%% A like-for-like comparison. The lemma statement is character-for-character
%% identical at the two checkpoints, so only the proof differs.
\begin{figure}[t]
\small
\textbf{Under \paco{}}
\begin{lstlisting}
Lemma eqit_inv_Tau b1 b2 t1 t2 :
  eqit RR b1 b2 (Tau t1) (Tau t2) -> eqit RR b1 b2 t1 t2.
Proof with eauto with itree.
  intros. punfold H. red in H. simpl in *.
  remember (TauF t1) as tt1. remember (TauF t2) as tt2.
  hinduction H before b2; intros; try discriminate.
  - inv Heqtt1. inv Heqtt2. pclearbot. eauto.
  - inv Heqtt1. inv H.
    + pclearbot. punfold REL. pstep. red. simpobs...
    + pstep. red. simpobs. econstructor; eauto.
      pstep_reverse. apply IHeqitF; eauto.
    + eauto with itree.
  - inv Heqtt2. inv H.
    + pclearbot. punfold REL. pstep. red. simpobs...
    + eauto with itree.
    + pstep. red. simpobs. econstructor; auto.
      pstep_reverse. apply IHeqitF; eauto.
Qed.
\end{lstlisting}

\vspace{0.4em}
\textbf{Under enhanced coinduction}
\begin{lstlisting}
Lemma eqit_inv_Tau t1 t2 :
  eqit RR b1 b2 (Tau t1) (Tau t2) -> eqit RR b1 b2 t1 t2.
Proof.
  intros.
  step in H; step.
  now apply eqitF_inv_Tau.
Qed.
\end{lstlisting}
\caption{The same lemma at the \cpaco{} and \cnew{} checkpoints. The statement is
identical; only the proof differs. The \paco{} proof interleaves coinductive
plumbing with an induction on the derivation, and six of its tactic calls
(\code{punfold}, \code{pclearbot}, \code{pstep}, \code{pstep_reverse}) exist only
to move between \paco{}'s guarded and unguarded forms. The port separates the
two. Taking one step in the hypothesis and one in the goal is \code{step in H;
step}, and the induction becomes \code{eqitF_inv_Tau}, an ordinary $19$-line
lemma about the one-step functor that contains no coinduction at all and is
reused elsewhere. The total is not much smaller. What changes is that the
coinductive reasoning is three lines and is separable from everything else.}
\label{fig:example}
\end{figure}
